\documentclass[12pt,a4paper]{article}
\usepackage[utf-8]{inputenc}
\usepackage[french]{babel}
\usepackage{amsmath}
\usepackage{amssymb}
\usepackage{geometry}
\geometry{left=2cm,right=2cm,top=2cm,bottom=2cm}

\title{Chapitre 4: Mathématiques}
\author{}
\date{}

\begin{document}

\maketitle

\section*{Contenu}

Chaqutre CE Toametures d'une eaqrace eucholoone.


Œ, A5) et au ev de olmentimu n, 
Lœ* encore 
T'Dametue - un mate (U U Li maœwxue adiacrte au gractect acaluses. 
D Defiatein. à On appelle trameture de E laut cuolurepliree 
Vx€eE Ice I - I (l 
Remarque : À ent um coaumarpehasine .

On a Le cornaclerusahans auivauter.


ancthms petvauntes 4œul eprivaleutes : @) Pere une coameGue. @) À camenne Le produit s&aloure: V3) E* PR fp <ex 1? G) Gramafarume me BON em re BOW 
(&) la anaburce de À dans tune BON at axthagauall TS QD 2 8) SG) A) 
Puuve : Hantbriaus Suppanans C4), an vappeUe qe fauve G,uy)e Ee° Cxla> = Ale ele ty) 
<£C1L&4)> - 4 CALE M 16) le IPGN) A (LL cru = LPEONT  (LPra d E) l'épuus 4) 1 ( Ÿ ‘é À Ftà TT 4 ( (rise ty) 2 Cxla > 7 Clerg 4 è 
dance 
ane ax à buem 7) => C2) 

e 2) 43) Suit L-fu.,out ane base axthanœamme cle E (ce démentenn n 26 DR folle jus, et au me mec fl) e lv <wlve> = € te) \ fees: <velUi7 = Se Cr iv = < fa)iPlop Cesce AT et une feslle ole n veeleurr archamarmee daurs € de dnmeurren n olaxe À ext ame Bow. 
«(37 (4%) En reprenant Les vataliwns Lu ren rt La aura Liu te u tu, de Solo An ton & de : Pate, 
H: ë eiÆ dise Lx amalute ole troasvage ele OT a F et aœvmme Les des boues sant ON ((3)), Le matruie oo acéhasauxle. QG) A). Saut ox, 00% une BON, x mec aLers .

z L'xsus £Æ 4e Je Z , que ales 
LE) H({)-nx 
an cour une RON, qu à ele xT x æ LEE VU Y 1x) x XTHTHX, Come D est exthaganle, MTH-T, et clame PO 2 ee LE let LP) le 
{l 
à 
n 
ec Cennuine D alémamnbie tin. Thaëréme : + Sat E um espace euchilien, À'euse Le oles coametres veelarreUles ole E } anate” OC) ex€Æ un ous gpreufre 
ele GLCE). . Su S ot me BON de E alars an à um caumarephénme 
le gear PR: OC) —ok) LL — He (@) 


& 2) Drapuecetes ol Came lever Vrepreéetes ol téoume bars 
x» Le determéuant d'une voameluie ol E vaut A ax - À Lie cumdtues de déterminant À Raul aggelers ARE» œux ceceeler eÆ ferme un gaie SocE) __ Le» cnemelaces le de -1 ant apquelees meakeves eu emdimeckes . CE N'EST PAS ON GROURE | , SC Fet un serv oe € sEalle gr Je OCE)D «laxs Fhamc: SO Veer, PRET als peur ge Come, Are un csamerphisme de E clans ET due Fe. Aires, D netrutlian de Ÿ Œ Fat au un coæmouuhcsinre de Fur F 
Main£imant  SaxG 4 € FH et + qgauelkaunque doms F3 Da A extite 2€ F7 E Plrrzx et olane 
La) x >= € A4) 1FP(z)5 = <u | zY7 -=o ?4 fw 8e 413 
Done Vue F7 Veerë, < f{)\=2-0 Va er” y) er. caro.


x & À a des malours joeges elles ans elles #aut clous 44. (Fous À parvereut- ua cle vap caumpleres.) 
x Les sev ex CP-TA Jet ken (PrTt4 saut arethogauause : Se D LE alara PA LR L) 4 
D Eten (xt) alou fly)=-4 ex La ly7= <fG)tPry) >= <æel-g>=-lyT 

3) Les Lgumetrsen ortheqauwelen 
Ragyel Sout F um sev de E lus FFE - 2 U et La Hymne athacçoale fe reapeprart ee F at La Hpmebrcce ec rapigesÆ ee F Jeellbemeert 
à F* 
On at 8e Pr pri: 2pe - Td LH -2pu 
Prequuiutes Les sygreebures ethaqunales Hant des isamelbuer En effet ja ger et 2€ FT ads lg + gl” (CET Set g *€ E , d'une fre 
æ < pee) 4pes Ce) olœue : LE (es il Pr CN + {l Pet Ce) Il et o'aubree gant 2 (&) = Pr) —Pri D: e Le > UT: PF + tt pres GI Te Hell.


2 
on deu Que HRMENÉE AE y dus Le cas les Afebrres 
NA Gun & kex (dk -Tl ) = F er er (os +, ) 2 F* dane À Dv2 olraganalisa lle et chu ume borne aclapetee 
Hatn if) pa PS 4 © attsut x À (4) o que La clés de F(F) On var'é iex lune bue ext olscecle 4» clin FTeré- qu'une Age ge (VOCABULAIRE : At laypequéux, he et une see flemae ) On œ uen Sûre bien el tudres CD œuvre CUS pars PA bueu Les clentefeve Ceu prardrie lier Les Aala Gens ) à fur haënr nue “auectatiqn “ ole L'espace. 

[s 
T Clsrc pion tcrun des (rame bus eu demarrer E.


Roume La efracrs euclolieus Le démeusinr Lj quicn saut laus aemarphes à IR, mn à en parkiiubier Le jhun mactauel qu'an ‘œurnte ! paurt £auarr cdamx quel veus camper Ars anses A) Sue_L'anentétian ot han.

eau E an ft Caunes um rmectlêur olans Le ua ) a camuentran usuel OL ol cannarleñcex qguee- Era le. mcerse  otes oucqulles l'ame nmentze : 
TS 3 
RTE CRÈE +T 5 FR)- = 20 RC CE,w)- Y 2% Cen] 
Was À Eet um hair anlutuoure mars esseutre(. pause Ctimpucemn dre ae er emyele comment agct tune hetaitièn ol ‘angle. &. Une fau alle canrwentiän Jexce , als une Las du plan (EF) eme dunecte a CET )2+T e£ Emelseeke se (P,#)=-E. 
Où marque qe fe exemple AC @&, P jet une BOuD als CP,T) et émotiscecle. 
Per Se (PT, T)= B el une BobD et (E,7) exbune BON alars 851,7) at edxecte SST Lx mature de 027 de de B'a B € olans SO (2). Sut ©: (È,©)) 
TD a > & =? un (1) Sat CE ,VIetE ee À =comOT + puu SA  V’=-MuD +30 
«{ d''œi Ps ce au à e SO). 
nE  <ce30 L , 
pres D can © © -BÊ <æn 
(rene ve) g Ste) 
ii 
eyfeut CUVE 

E Plus qencraliment , dus un emace euclolen cle dementéen n , œum +e olarme une BON Le clenaxtBget au à ) one 8 chan cœummne etant celte par contents et elases 
Pet me BOND n°  Har(D Es) € SO).


«Planta. le cha cle Bone fut qar Pirates 
e Vans R?, an Barrcé qe DA Lare carnage. sou tolieele 
> Dans E qebeeuque 0 faut Pever umue cameeriteam qu qu” Etre culutrnise. 
+ Dans À'enpoce Ce IR?) mœu revemcolrauns qrhes Guol sur Ja cmuentian vue.


2) Docurliau de Ole). Thedëme: SU P- (a) € OÙ) aber 
Saxe r-(à x) mue attb 4 (eSO() ) Sat €, (° b } aa (CE DE)\S0()) b 74 
Preue. æ dub A e& c°+l°: À (a,b)= (cam 6,546) Ce, el): Ceana, ua) eÆ comme La muakrure. e£ L ) Orac+tbel = co ce + aën © Heu = con (4-8) A ae pot + 4-62*+{ cerf cas à ses (S+E )- 4e 6 em à = mu (@tT) = ca @ . 4-0: -T rer] cab à >: coè(e-T) = Hu 
Men à = samÛB-E )- -4280.

cart. 

Π
2) Carsfieu ae en démensieu 2. Th, Su € un wpace euchlier ol olemeutier Ze À ume BOUD. Sœut PEOC), CRE - Ain € 
1) $esoce) æ 346 bel que Hat (A) « H 
AÆtne cas © te) : Dans ce cas et une ralakun d'au & e£ 4x 
EE ee Ÿ4 -TA, $ m'a por cle vap sceelles.


2) Le ENS) > T6 tel que Mats(f): (Le eo) Daus ce cas mL une eflestius fee repart La.

ébaite D: Lex c- Tl, ) et A existe uue BOND Lun 
L dupe We À #C gaur mature 52).


Preuue ; On raqace Île ge Le déterminant Alan enclaumerghisme e € ciholepewolant cle La bare chaire . D PESOEY ae Abe def (rat ?) = 145 HeSOG) Le 
M [* Am © [- K7_2Xca@ + À À: 4 Ces -4) #lane ÿ d=o nu 6-6 [TT Æsemau Ace, 2) Où à L'épuuvaleuce arr Da meme Haisan. b 726 rene = x2- 1 elane Sp €) = j4,-45 
— sim X + cesse 
- An © x 06 
fe oragonalisalle Sert D umutsure de er (P-T4 ) PE ke (Pre) Aus DIV Hate) À (2 °) 
& A 
am nécamatt Lrou x La Age Bee L ae 77 Lex (P_ Ta, ), Le 
es des ralpthss 


[8 
3) Praprcetes des fa Chine . 
e Va,f Ha Ha GG z asp. 
« On em leduct qe BA Ÿe SOCE), 52: satrsd Ha nrË bEpert À > “LL °% Bowb ehaire. EE olane € (mao ln 2T) me olefceucl fers cle Lx Bon? chounce . 
se PB e£ D! name 2 BOND et P- Hat (8,2) 
AC FM: Faks eX m2 Hal, alu .


Me PMP 2e Mer: H.

Hp Mo Ma 
Où male als kams amdegec Éc' +- Fe .


Prapuete: Se Mu lE)- 2 SO(E) ex# Le groupe cles natalkiimn qu est commmctalef. Cruorpsrpe res Jus F2: ri ex) 
4) Detemmcnalion le caraeberusErigecen cime coame tree ex? dinanes 2 : 
SE œu caunoeut La menlaeree Mol une inameture dans ume BOAD eéreete alaré. an feinnce. 
- et A1 PA 7 H=-4.

HET MT He La draute Dikex (P-xd, ) =Td, Ÿ: au determine Led ge de Le atalism . ÿ- b5. 
La atuntinn our Les capaces de démeuran 3 44 un 
peu quan camplexe. 

D D man olémeubssr 2.


D Dites et fractuu E vec£arcæel 
Lottease : cÀ fur aetiden que bases Aaut” cliecles .

Dans À 'erprace., am 4e fixe Lun. qurem ces vecteur €, e£ ane Cholilt nu 4e peut) cle Halakèan aulaur de e, À. er L'er À regle le LE borelez =. Lrepac el au laurrre ke.

Caen enr. ler lier Vanelansanmens us. rentes Une fac elle vésaeekess else ee sépare" Les Bow on BON directes et ieolireelen, 
Arr carter kr ea, 2,3 ) GND.


Ex SE (ex, er) Es ) Directe Ce , CRE) chrolsnecÀe 
T8 le, se, POLE A 
BARS + 
Daru RTE casa late quee- LG bue camrnanne quee- et uuecte. Stan ot ol eu ace ele ein Lau 7% TRE 9 Le. ehacx ane Re DORE RON) 
Praprane tram 2 définit SuEt E un egace eucl'dess axveute 
NE CR cle €. feux Eau x4,2, le € Vuce delg Cent e) 2 CWIR7.


Cm appelle w Le grace vedtæxeR de Lt, W=2Æ,A #2.


da ef dt a na 2 … x $) Dance Das 8 Eu L 7 œ peur caen cle ,mmees (E) Le SA LOS KA <e pet d 
ec x= (Ë) at Wd:.d: à 
Ra Per Ê 
PR (BK - ë 5) A — B (2-28) + (APE chetg Pate Xs 7 | ml 
du à Burns wo Va PES Bite - ape bre WW - (ER) red 
da Pe dal 

Remmenque À % Se B'ekt une cure ÆCND Ares Lx eo 
mhaleise cle gesege ce BP 'a À Pl ext daus SOUS ) clans D x, — x! = PX4 NÉ PX) KE PX et se We pP\ alas oleÆ. ARTE TR: 2) REP A pu) 
1 
ee <W'X>e #puyt Px', WTP'px 2 w'x dire Le mecteuwr WW me clejeunel {6 le Un Beub Here Remarque Li Pau 4e pouuencr de L frurule 
a À Wa Les ra Lu 
Es Prapucetes a) CAR EE TE eÆ peur AE lR Spa tx A LS = NMecazs) à we) (Au£inumetuce) 24 À (a kpzs) = À Gil) flex, (Belimrireite) () rina,-0 > met R, Aout coleneniees. (3) Sc mA +0 alans xn% et euthaganal à Vect{ a, x,? et A Are réa, femelle obhamace aa Îæ Le, AAA %2G A me @BovD eliecte.


I) &E une cammequeunce de Le miltiliseate" dx déterminant et de fat que olebCes, 3,2) 2 + olete (es +) 
(2) Se œ,,0, colémenire : Vx CE Lt Ce, ,x)=2O = Caux, le > donc AAAN, = © eË 2x +, A%27:0 , Ve ce lrE£ (ris je) < ©. Se ei, m1 leloueut 4 clseoures ader am passeur carmpele Ex en une bue de E Â2,,x,,æ3% mous okaus clef (42. ,æ3) >+o0 | &) far oclleuro Re A >= Rite, k, #4) = D + EËT ele mudime aure Xe. ele lt = VecbS aa ze &s 
r$: Bo 
Son y nr aflans cel, (x,x Lu Let (ent za: 
4 Wear 8 ! D 2,4) rs Ga 2} TA 2) A ° = Maur il 
ane. ni %2, 4 ) EE OND et rare À À 


4 Obs 
Ramanqgee : Se Cove, 03) est une BOND 
“ VRAU3 Z Ua A0 3 E AUS ee un _, dun À 'epace ï DEVEE a “: U:. y Ée —s dans ua BOND Cou. vs te exrwpele D Uy À V3 “4 7e] sl [2 b el LA ©. A OV” 
Pradut vetanseN 2€ au RE.

ee tuyU2 À vecteurs man celimensres ef Fe Vect 15. pa On. chainit me Bow cle gham m Æ,r) ee Le. que va = lu © 
ul (cmeT tam Se LP) ne 2 Ag albses 
re 
CE 
À z| _ el S = Qullllu ll ces & 
FRE _ 
DER an Us = lol Aa © e.


2 
+ 

29 ann Pre a rte £ee PS Te 3 
Three : Sat. El nue eafrace. be Le ben: = elceute" | velarecele . 
A) Se À ÉSOCE) et one: chaine Gcce ertire., d ere un nel & et Bne BonwD ES Le E Eh 
Le \ (eu z, 3) 
_ X | ex + dune case Gere 
À S & Haks (?) - [: CRDHÔ An © 
CO ACUË ces © 
Pa ee de Ÿ PE Le natal d'une D due pere, 
et d'a & ; D es ytea) | S'é-0 A et ; De v(es D À de ar Bens. olineetr. ) 
RNA © 
ce _ 
29 S> PE CUNSOCE) bars LL exrnte un see @ et une BOUD (ee, e:) &B qe. 
LA © eo FRE 5 (4) = Ë case ue ) D One ce E& Ds ce cas DONS ER at campus ee de - L hatatins H'ane © ohnurpe arc eu eÆ olhayle s La rflencrsr. fes rapqarte à D+- Vest Se.,e.f 
. Se &:6 abrs 
À 7.8, 
Ua) 
Maur nevceucluams Au D foune js Eanol.



Êe Le € mé um enprace eucllien , de A ane BP de € et J une coamelrte. dan rx cernes EC PA inatecelllss Le Le 4, or Céumineut rater. 4en cerclreers lcqacen ?

+ las Ans SOU ME Sor2) ; als anrart qui ‘À eointe Peso) 
-A . PTHe - A © © ba PTHP,: PHP (: Lt me) dams une pose le,,2,e:) 
8 An case eo Su Mz T, eÆ alars Ÿ= 3e : . Sut. My JL s un determine. LP ae De ken (Ÿ- il, ) e£Æ em olous<t an Veclerue er ele marie A Ca DE teet{ef x Comme Lx Hrance ef cuuescaute- 1e Sutole base, Te (1) = Mc Cine ar cenmmait @ au Age F2. (ee C-mtr1) Lx net par alraure à €, (xsue,+Be+Fe ae teum, tx) < A & à s o 8 pce -5 ue Oo X fan e 4Ycare N Cats, Bruce Cer,e,,e3) BOND au ohesit &> acthagsl a ot, ef am pare Er SE Ae, . Ceres CE ea £ akhnn une BOND £Æ clans cette bare Mirrehessn ele fre altemolue , 2, es _. 4 © (a \ , Cas 1 M € C(3)\$C() TE cussle TC £ ë P HP: (ea - 40 en D CRT s Sa Ÿ= - (ta: mi) s Sat D 2: kex {8 + TA, Ç = kex {nm+x) De ete FUME: Arr e ROUE CT TT eÆ ohmauneat, AC 2€ D 2e£ (ex, Pt=)) ol Here le 25e, + On Pat came class D2 ces À. 
3 (Pr) m0. 

(43 4) Guuue du écéme de Lonhintinn. Î 
La gone. ot esresme à e far une alrernnfriia… 
Mere Les grluanmen caracleresqeees .

Set LEE), e(xrA-f) 2 X'4a KL X - detf.


Cammre- eg E un 4 me oe degree 5, «À à part 
ume seule nacme. se: » Bart 3 
Canme À af un crameluie 
êbre que À Æ fau A. 
y Les raanes retllen me peuueatr 
seen OŒ)\SO(E ) X69- {(x-4)° X/ Ge) 2044)?

Xyex bay) 3e (XI = (X0) G- 2) ° Hg (x )= 0-4) 0c-p) Xp (X) = (KA) CX- BI -R ) 3 ix=f) À re ne et P = — ipl Te def Ÿ: 1BI= A 
A0 'Se AE SC(ET an mut que À erÆ elles  . le d.

dame an ehounxt €, ES qe e œu 
Pées > = e FaVectiest canme F et stable par À, Faure, et a nebuction de À à FT aé une csamelrrer fe, E FT urbaine eÆ e;se,1e, alars À meobluce cote À dans La Hone fées 4 lecsat feu <a (es) 
A © © 4 P- Fabio, et (4) 5 | e 
© a b Le 
ee el 
une qarË en rmenttcehisnt ot Je, ,e,% ; an a une cramebeie ne (É£)e go) et À = let P= elet (É 5 
e 
Aaue c'at une nalitinn ox han: «À exrls 6 
bel que 41e P= eee tp? le ce - 440 
© AëuË CRE 

ET = S ŸE St)\sote) ) “A1 GE valeur frape le ee À. Camnrme Are den eur ur chratt - on vecteur Jraeqre unibaire €, GEfteraer) jy um gare F2 Veetfe.f LE an quencde €, nilaire oblams FT 4 Es = En ee ee a 
41 ee le] Fe FE es + … ' e Ca : (2) < æ(ix)e OU) mes cn clac ue Je 
ut L'enioteuce de ER El qe d e ce 
A Mob es ei É- EC DE ue ©  Lenû ces CFD 
—_——— 

\end{document}
